%% kapa: exam=1 round=36 year=2025 period=2 form=A subject=accounting subject_ko=회계학 seq=4 number=44
%% source: Q-Net 감정평가사 기출문제 (cst00309) artlSeq=5231525 file=raw/1008/5231525/2025년 제36회 감정평가사 1차시험 2교시 문제지 원본.pdf
%% conversion: kordoc 4.24.51 -> tools/gen1.py (auto); review=proofread-sonnet; flags=none
\begin{question}{44}
(주)감평의 20×1년 기초상품은 ￦20,000, 당기상품매입액은 ￦80,000이다. 20×1년 기말 상품 관련 자료는 다음과 같다.
\begin{tcolorbox}[enhanced, colback=white, colframe=black, boxrule=0.4pt, arc=0pt, left=4pt, right=4pt, top=4pt, bottom=4pt]
\begin{tabular}{@{}l r@{\hspace{3em}}l r@{}}
○ 장부상 수량 & 90개 & ○ 단위당 취득원가 & ￦300 \\
○ 실제 수량 & 80개 & ○ 단위당 순실현가능가치 & ￦270
\end{tabular}
\end{tcolorbox}
감모수량 중 7개는 정상적인 것이며, 나머지는 비정상적인 것이다. (주)감평의 20×1년 포괄손익계산서에 표시될 매출원가는? (단, (주)감평은 정상적인 감모손실과 평가손실은 매출원가로 처리하고, 기초 재고자산평가충당금은 없다.)
\begin{choices}
  \item ￦75,100
  \item ￦75,400
  \item ￦77,500
  \item ￦78,400
  \item ￦79,300
\end{choices}
\end{question}
\answer{3}
